\documentclass[11pt]{article}
\usepackage[a4paper,margin=27mm]{geometry}
\usepackage{amsmath,amssymb,amsthm,mathtools,booktabs,hyperref}
\hypersetup{hidelinks}
\newtheorem{theorem}{Theorem}
\newtheorem{lemma}{Lemma}
\newtheorem{proposition}{Proposition}
\newcommand{\R}{\mathbb R}
\newcommand{\norm}[1]{\left\lVert#1\right\rVert}
\newcommand{\ip}[2]{\left\langle#1,#2\right\rangle}
\title{A computer-assisted closed free-elastic expander\newline with escaping arclength centroid}
\author{Computer-assisted proof with two independent reviews}
\date{11 November 2026}
\setlength{\emergencystretch}{5em}
\begin{document}
\maketitle

\noindent\textbf{Status.} All interval computations described below have
completed successfully. The proof combines these computations with ordinary
analysis; it is not a full Lean formalization. Two independent analytic/code reviews have completed successfully.
End-to-end Lean formalization of the ODE, analytic estimates, and geometry
is not claimed. The complete exact decimal
center and interval output are in \nolinkurl{certificate_center.json}
and in \nolinkurl{validated_certificate.json}; the reproducible implementation is
\nolinkurl{arb_shoot.py} with \nolinkurl{validate_certificate.py}. The latter
records SHA--256 hashes of the former and of itself.

\begin{theorem}\label{thm:main}
There is a smooth immersed closed
free-elastic flow whose arclength centroid diverges to infinity.
Its initial arclength centroid can be chosen to be zero.
\end{theorem}

Here the arclength centroid of a closed immersed curve is
\[
 c_\gamma=\frac{1}{\mathcal L(\gamma)}\int_\gamma\gamma\,ds.
\]
The construction concerns smooth immersions and allows self-intersections.
It does not establish the migration property.

\section{Polynomial shooting problem and geometry}
For the state $z=(x,y,u,w,k,v)$ consider the autonomous polynomial system
\begin{equation}\label{eq:ode}
 x'=u,\qquad y'=w,\qquad u'=-kw,\qquad w'=ku,\qquad
 k'=v,\qquad v'=-\frac12 k^3+xw-yu.
\end{equation}
For $p=(a,b,L)$, let $z_p$ have initial condition
\begin{equation}\label{eq:initial}
 z_p(0)=(a,0,0,1,b,0),
\end{equation}
and define, whenever this solution exists up to $L$,
\begin{equation}\label{eq:shooting}
 F(p)=(y_p(L),u_p(L),v_p(L)).
\end{equation}
We also integrate $m'=x$, $m(0)=0$; this seventh coordinate does not affect
the first six equations.

\begin{lemma}\label{lem:reflection}
If $L>0$ and $F(a,b,L)=0$, then the positional part of $z_p$ extends to a
smooth unit-speed immersed closed curve $\beta$ of period $2L$ satisfying
\begin{equation}\label{eq:soliton}
 k_{ss}+\frac12 k^3=-\ip{\beta}{\nu},\qquad
 \nu=(-w,u).
\end{equation}
Its arclength centroid is $(m(L)/L,0)$.
\end{lemma}
\begin{proof}
Writing the six-dimensional right-hand side as $f$, the linear involution
\[
 R(x,y,u,w,k,v)=(x,-y,-u,w,k,-v)
\]
satisfies $f(Rz)=-Rf(z)$ by direct substitution. Both $z_p(0)$ and
$z_p(L)$ belong to $\operatorname{Fix}R$. Thus
$Rz_p(2L-s)$ on $[L,2L]$ solves the same ODE and agrees at $L$ with
$z_p$. Local uniqueness identifies this reflected solution with the
smooth continuation at the seam. At $2L$ its value is
$Rz_p(0)=z_p(0)$, so uniqueness also joins the endpoints smoothly and
gives a periodic solution for all real $s$.

The identity $(u^2+w^2)'=0$, together with the initial data, gives
$u^2+w^2=1$. Therefore $\beta=(x,y)$ is a unit-speed immersion,
$T=(u,w)$ is its tangent, and $T'=k\nu$. The final equation in
\eqref{eq:ode} is exactly \eqref{eq:soliton}. Reflection gives
$x(2L-s)=x(s)$ and $y(2L-s)=-y(s)$, hence
\[
 \int_0^{2L}\beta(s)\,ds=(2m(L),0).
\]
Division by the length $2L$ proves the centroid formula.
\end{proof}

\begin{lemma}\label{lem:flow}
If a smooth closed unit-speed curve satisfies \eqref{eq:soliton}, then its
dilates $\rho(t)\beta$, where $\rho(t)=(1+4t)^{1/4}$, evolve by the
normal velocity of free elastic flow. A smooth time-dependent
reparameterization yields a solution of the strict vector equation
\[
 \partial_t\gamma=-\left(k_{ss}+\tfrac12k^3\right)\nu.
\]
The centroid of this solution is $\rho(t)c_\beta$.
\end{lemma}
\begin{proof}
Under dilation by $\rho$, the scalar $k_{ss}+k^3/2$ is multiplied by
$\rho^{-3}$. Since $\rho'=\rho^{-3}$, the normal component of
$\partial_t(\rho\beta)$ equals
$\rho^{-3}\ip{\beta}{\nu}=-\rho^{-3}(k_{ss}+k^3/2)$.
For the vector equation let $\psi$ solve on the parameter circle
\[
 \partial_t\psi=-\frac{\rho'}{\rho}
    \ip{\beta(\psi)}{T(\psi)},\qquad \psi(s,0)=s.
\]
Its smooth periodic vector field gives a smooth circle diffeomorphism
for every finite $t\ge0$. In $\gamma(s,t)=\rho(t)\beta(\psi(s,t))$
the tangential velocity cancels. Centroid is unchanged by
reparameterization and multiplied by $\rho$ under dilation.
\end{proof}

\section{Uniform analytic estimates}
For the augmented seven-dimensional state set
\[
 W=(1,1,4,4,16,64,1),\qquad
 \norm{z}_W=\max_i\frac{|z_i|}{W_i}.
\]
All operator norms below are induced by this norm; multilinear operator
norms use it on every argument. Consider the convex tube
\begin{equation}\label{eq:tube}
 |x|\le1.6,\quad |y|\le2.4,\quad |u|,|w|\le1.1,
 \quad |k|\le3.5,\quad |v|\le6.
\end{equation}
There is no restriction on $m$, since the field is independent of $m$.
Absolute row sums of the polynomial coefficients yield throughout this tube
\begin{equation}\label{eq:derivatives}
 \norm{f}_W\le2,\quad \norm{Df}_W\le8,\quad
 \norm{D^2f}_W\le44,\quad \norm{D^3f}_W\le192.
\end{equation}
For explicit verification, the oscillator Jacobian rows are bounded by
$|k|+4\max(|u|,|w|)\le7.9$, and the final row is bounded by
\[
 \frac{|w|+|u|+4|x|+4|y|+24k^2}{64}<8.
\]
The $k'=v$ row has norm $4$; the remaining linear rows have norm at most
$4$. For the Hessian, the oscillator rows have norm $32$, while the
final row has norm $12|k|+1/4\le42.25$. The sole third derivative is
the cubic term in $k$, and its weighted norm is $192$.
The field norm is less than $2$; the largest simple bound needed is
$|m'|=|x|\le1.6$.

\begin{lemma}\label{lem:taylor}
At any real center $C$ in \eqref{eq:tube}, the exact solution of the
augmented ODE starting at $C$ has a convergent Taylor series for
$|h|<1/12$. Its displacement is at most
\[
 w(|h|)=\frac{(1-12|h|)^{-1/2}-1}{3}
\]
in weighted norm. If $P_N(h,C)$ is the degree-$N$ Taylor polynomial, then
\begin{equation}\label{eq:tail}
 \norm{\Phi_h(C)-P_N(h,C)}_W
 \le\tau_N(|h|):=\frac{(12|h|)^{N+1}}{3(1-12|h|)}.
\end{equation}
\end{lemma}
\begin{proof}
At $C$, the absolute coefficients of the polynomial expansion of
$f(C+q)$ in the increment are bounded by the scalar polynomial
\[
 2+8q+22q^2+32q^3\ \preceq\ 2(1+3q)^3,
\]
where $\preceq$ means coefficientwise domination. These bounds are
absolute row-sum bounds, so they dominate the convolution recurrence for
the absolute values of all Taylor coefficients of the solution. The scalar
majorant $w'=2(1+3w)^3$, $w(0)=0$, has the displayed solution. Induction
in the coefficient recurrence proves coefficientwise domination. The
coefficient of $h^n$ in $w$, for $n\ge1$, is
$12^n\binom{2n}{n}/(3\cdot4^n)\le12^n/3$.
Thus the vector series converges absolutely for $|h|<1/12$, solves the
polynomial ODE by termwise differentiation, and agrees with its unique
local solution. Summing the coefficient majorant from $N+1$ onward proves
\eqref{eq:tail}. This argument uses the derivatives at $C$ only; it does
not assume beforehand that the entire local solution remains in the tube.
\end{proof}

\section{Validated integration and its interpretation}
Let $p_0$ be the exact decimal rational triple in
\nolinkurl{certificate_center.json}, approximately
\[
 (1.36288130730646586725,\quad
  3.01173359847420887094,\quad
  7.49497740159149112339).
\]
The finite-difference step is $\delta=10^{-70}$ and the desired fixed-point
radius is $r=10^{-80}$. Every integration uses $N=160$, $n=750$,
and $h=L/750<0.01$. The exact, unrounded Taylor iteration is
$C_{j+1}=P_{160}(h,C_j)$ with the specified exact rational initial state.
Arb ball operations enclose this entire discrete iteration, including
the input conversion, coefficient recurrence, and polynomial evaluation.
No midpoint reset occurs between integration steps. Hence every exact
$C_j$ belongs to its computed ball, and the final ball already includes
all propagated arithmetic error.

Every computed grid ball is verified to satisfy the smaller bounds
\begin{equation}\label{eq:inner}
 |x|<1.45,\quad |y|<2.26,\quad |u|,|w|<1.001,
 \quad |k|<3.02,\quad |v|<4.15.
\end{equation}
The smallest weighted margin from these bounds to \eqref{eq:tube} is
$0.02475$. The verified inequalities include
\begin{equation}\label{eq:globalerror}
 w(0.01)<0.023,\qquad
 750e^{60}\tau_{160}(0.01)
 <1.816799\cdot10^{-120}<\varepsilon:=10^{-118}.
\end{equation}
In particular more than $0.001$ of weighted tube margin remains after
subtracting both $w(0.01)$ and $\varepsilon$.

Here is the precise comparison with the actual ODE, which is essential:
by Lemma~\ref{lem:taylor}, each local exact arc $\Phi_s(C_j)$,
$0\le s\le h$, lies in the outer tube. As long as the actual solution
also remains in that tube, its difference from this reference arc is
bounded by the flow Lipschitz factor $e^{8s}$. At grid times the errors
therefore obey
\[
 e_{j+1}\le e^{8h}e_j+\tau_{160}(h),\qquad e_0=0.
\]
At grid times and between them this gives a bound at most
$750e^{60}\tau_{160}(0.01)<\varepsilon$. The spare $0.001$ margin
prevents a first exit. Standard continuation for a polynomial ODE then
gives existence throughout $[0,L]$: the first six coordinates stay in a
compact set, while $m'=x$ bounds the seventh coordinate on this finite
interval. Thus adding the ball $[-W_i\varepsilon,W_i\varepsilon]$ to
each raw final state component gives a true ODE enclosure.
Arithmetic radii must not be multiplied by an additional Gronwall factor,
because those radii already propagate through all discrete steps.

The exact center also satisfies $0<L_0-\delta-r$ and
$L_0+\delta+r<7.5$, as verified separately by Arb.
The same argument supplies the needed uniform neighborhood. On the box
\[
 \norm{p-p_0}_\infty\le\delta+r
\]
the initial-data perturbation has weighted norm at most $\delta+r$.
With $E=e^{60}<2\cdot10^{26}$, the change in state on common time
intervals is at most $E(\delta+r)$. A possible extra endpoint interval
has length at most $\delta+r$ and costs at most $2(\delta+r)$ in state
norm. The computations explicitly verify that
\[
 (E+2)(\delta+r)+\varepsilon+w(0.01)
\]
is less than each weighted inner-to-outer margin, with a remaining margin
greater than $0.001$. The first-exit argument, including the possible
short endpoint-time extension, therefore places every intermediate
parameter trajectory in the same tube. Consequently the derivative bounds
below hold throughout a convex parameter neighborhood containing all
finite-difference segments and the fixed-point cube, not merely at the
seven evaluated parameters.

\section{Derivative estimates and the fixed point}
The state-flow variational equations and \eqref{eq:derivatives} give, for
$0\le s\le7.5$ on these trajectories,
\[
 \norm{D_z\Phi_s}\le E,\qquad
 \norm{D_z^2\Phi_s}\le\frac{44}{8}E(E-1)<6E^2.
\]
The latter follows by integrating
$44e^{8(s-t)}e^{16t}$ over $0\le t\le s$.
For derivatives involving the terminal time,
\[
 \norm{\partial_s\Phi_s}\le2,\qquad
 \norm{\partial_sD_z\Phi_s}\le8E,
 \qquad \norm{\partial_s^2\Phi_s}\le16.
\]
The initial-state map $(a,b)\mapsto(a,0,0,1,b,0,0)$ has operator norm at
most one from the parameter maximum norm to the weighted state norm.
Extraction of $(y,u,v)$ has norm at most $64$. Thus on the whole parameter
neighborhood the shooting Hessian satisfies
\begin{equation}\label{eq:H}
 \norm{D^2F}_{\infty\to\infty}
 \le64(6E^2+16E+16)<H:=10^{56},
\end{equation}
where the left is the bilinear operator norm.

For $j=1,2,3$, compute true endpoint enclosures at $p_0\pm\delta e_j$.
The interval matrix $J_{\rm fd}$ encloses their exact centered difference
quotients. For each entry, the fundamental theorem of calculus and
\eqref{eq:H} give
\begin{equation}\label{eq:FD}
 \left|(DF(p_0))_{ij}
 -\frac{F_i(p_0+\delta e_j)-F_i(p_0-\delta e_j)}{2\delta}\right|
 \le\frac{H\delta}{2}=5\cdot10^{-15}.
\end{equation}
The computed intervals already include endpoint evaluation errors;
these errors are not added a second time.

Use the exact rational matrix whose entries are the following finite decimals:
\[
 B=\begin{pmatrix}
 -.06037342357664187&-.27213682312533344&-.05965645081977437\\
 .013384830972815771&-.0024448406036154074&.061329210847373104\\
 -1.0100953172650646&.053420898617877965&-.08577738110346352
 \end{pmatrix}.
\]
It is nonsingular. With $d=\norm{I-BJ_{\rm fd}}_\infty$, all matrices
in the interval enclosure are understood in this upper-bound calculation.
On the closed cube $Q=\{p:\norm{p-p_0}_\infty\le r\}$, the map
$G(p)=p-BF(p)$ has derivative bound
\begin{equation}\label{eq:q}
 \norm{DG(p)}_\infty\le
 q:=d+\norm{B}_\infty(3H\delta/2+Hr).
\end{equation}
The factor $3$ converts the entrywise finite-difference error to a matrix
row-sum bound. The $Hr$ term already uses the bilinear operator norm in
\eqref{eq:H}, so it needs no additional dimension factor.

\begin{proposition}\label{prop:computation}
The reproducible Arb computation, using 180 decimal digits, verifies all
grid, tube, tail, and derivative inequalities specified above, and
\begin{align*}
 \norm{B}_\infty&<1.15,\qquad |\det B|>0,\\
 d&<9.957524\cdot10^{-14},\\
 q&<1.168147\cdot10^{-13},\\
 \eta:=\norm{BF(p_0)}_\infty&<6.74\cdot10^{-118},\\
 \eta+qr&<1.168147\cdot10^{-93}<10^{-80}=r.
\end{align*}
All seven integrations passed their grid checks. The raw endpoint
rounding radii are of order $10^{-172}$, smaller than the explicit
truncation budget. The computed quantities are enclosed by Arb balls;
all acceptance comparisons use ball inequalities against rational
thresholds, not floating-point comparisons.
\end{proposition}

By \eqref{eq:q} and Proposition~\ref{prop:computation}, $G$ is a strict
contraction of the complete cube $Q$ into itself. Banach's fixed-point
theorem gives $p_*\in Q$ with $G(p_*)=p_*$. Nonsingularity of $B$ implies
$F(p_*)=0$. This proves exact closure, rather than approximate numerical
closure.

\section{Certified centroid and escape}
The weighted state change from $p_0$ to $p_*$ at their respective endpoint
times is at most $(E+2)r$. Including the moment coordinate in the integration
therefore proves
\begin{align*}
 m_{p_*}(L_*)&\in
 [3.78365096320749419634974942366049493295828418426801828
       \ \pm\ 2.44\cdot10^{-54}],\\
 \frac{m_{p_*}(L_*)}{L_*}&\in
 [0.504824866103541540109480685758874180664870977318606015
       \ \pm\ 4.41\cdot10^{-55}].
\end{align*}
The code directly checks $m_{p_*}(L_*)>3.7$, $0<L_*<7.5$, and
$m_{p_*}(L_*)/L_*>0.49$. Lemma~\ref{lem:reflection} produces a smooth
closed expander with nonzero centroid. Lemma~\ref{lem:flow} now gives
$c_{\gamma(t)}=\rho(t)c_\beta$, which diverges because $\rho(t)\to\infty$.
Translating the entire flow by the constant vector $-c_\beta$ preserves
the equation and yields initial centroid zero and later centroid
$(\rho(t)-1)c_\beta$. Together with the verified computation and analytic estimates above,
this proves Theorem~\ref{thm:main}.

Finally, the same intervals show $a_*>1.3$ and $x_{p_*}(L_*)<-0.22$.
Both corresponding profile points have $y=0$. Their connecting segment
contains the expansion center $0$, so $0$ belongs to the convex hull of
the profile. This example therefore supplies no half-plane separation
from the expansion center and no homothetic migration conclusion.

\appendix
\section{Exact rational center}
The following finite decimals define the exact rational center $p_0$.
For each coordinate, concatenate its displayed lines without whitespace.
This is the same center read by the validated computation.
\begin{verbatim}
a0 =
1.36288130730646586724540492564119328606333202087402037713638055
2888014387264453369650391025834236861156205915573992900599640530
66281049604343152049694198272133019826995101310866984

b0 =
3.01173359847420887094497074533823205110835009467373153380441310
3791670594510452679593380852864667621400092821914670389887258596
00926748645647853176928567838285224785930939511967857

L0 =
7.49497740159149112338537760573652397915893177644288585501917218
5315246830265113994711601095140425363807428194721693050936666465
29738270751311706827855707016904769809511819177453755

\end{verbatim}

\end{document}
